\scanpage{16}\mspage{13}

appliquant $\mathbb L_{\unsure{0|\circ}}$ en $L$, \uline{et les} $\mathbb{\pi}_i$ \uline{en les} $\pi_i$ \deleted{\unsure{définis}} \deleted{\unsure{satisf} la graduation de $\unsure{E|\mathcal E}$ par (1.3).} (donc $\pi_i\in E$)

\margin{condition essentielle !}
% slanted
\margin{\deleted{(iii.) Il existe \deleted{\unsure{une représentation}} un \unsure{$E$-module} \uline{fidèle} $M^*$ de $E^*$, \deleted{\unsure{fiel}} que les conditions du th\ldots\ du n\textsuperscript{o} 6 soient satisfaites (pour la graduation envisagée) (iv) pour \uline{tt} module fidèle \ldots}}

De plus, sous les conditions, l'hom annoncé dans (ii) est unique, \deleted{appliquant} et avec les notations du n\textsuperscript{o} 6, il transforme $\mathbb{\Lambda}_0$ en $\Lambda$, les \deleted{\unsure{$\pi_\alpha$}} $\mathbb{\pi}_{i\alpha}$ en $\pi_{i\alpha}$, $\mathbb{\pi}_i$ en $\pi_i$, \deleted{\illegible{2}} \inserted{$\varphi_i^{\beta,\alpha}$ en $\varphi_i^{\beta,\alpha}$}.

\uline{Dém} (i) $\Rightarrow$ (ii) Résulte de la table de multiplication des $\varphi_i^{\beta,\alpha}$ donnée au n\textsuperscript{o} 6, et dudit th. donné au n\textsuperscript{o} 6.

(ii) $\Rightarrow$ (i) \deleted{Soient $\pi'_i$ les images des $\mathbb{\pi}_i$. On aura $L^\alpha\pi'_i=\pi'_{i+\alpha}L^\alpha$ \unsure{voit|fait} alors que les $\pi'_i$ définissent la graduation de $E$, puisque ils sont \unsure{de|du} degré 0. Je dis que $\pi'_i=\pi_i$ pour tt $i$. Cf Proposition ci-dessous pour ce point, appliquée à $\Phi\unsure{-|\to}\operatorname{End}(M^*)$.}

Comme on a, pour tt $i\le n$
\[
 (\Lambda_0^{n-i}\mathbb L_0^{n-i}-1)\pi_i=0,\quad
 (\mathbb L_0^{n-i}\Lambda_0^{n-i}-1)\pi_{2n-i}=0
\]
on en conclut, appliquant $\varphi$ (et posant $\Lambda'=\varphi(\mathbb{\Lambda})$)
\[
 ({\Lambda'}^{n-i}L^{n-i}-1)\pi_i=0,\quad (L^{n-i}{\Lambda'}^{n-i}-1)\pi_{2n-i}=0
\]
Donc la condition (ii) du th. du n\textsuperscript{o} 6 est vérifiée, donc aussi la condition (i).

Pour l'unicité de $\varphi$, \deleted{\unsure{et donc sa proposition en}} \deleted{\unsure{et les autres}} il suffit de voir que $\varphi(\mathbb{\Lambda})=\Lambda$, et \inserted{pour avoir $\varphi(\mathbb{\pi}_{i,\alpha})=\pi_{i,\alpha}$,} \deleted{\unsure{au n\textsuperscript{o} 2 en réalité,}} compte tenu de $\varphi(\mathbb{\pi}_i)=\pi_i$, \deleted{\unsure{puisque des} \illegible{1}. Pour $\varphi(\mathbb{\Lambda}$} \deleted{\unsure{du fait}} \inserted{on utilise} la caractérisation \inserted{simultanée} de $\Lambda$ \inserted{et des $\pi_{i,\alpha}$} par les formules (6.\unsure{14|15}).
% archive sheet number: 15
